\section{Open problems and method}
\label{sec:open}

\subsection{Open problems}

\begin{enumerate}[leftmargin=1.6em]
\item \textbf{The general-$n$ tower from first principles.} Prove the Dickson
  factorisation (Theorem~\ref{thm:tower}) for all $n$ by exhibiting the functorial
  $\Sym(W)$ construction of \eqref{eq:symw} --- equivalently, assemble the bounded
  rank-one $e_2/\Lambda^2$ closure into the exact trace map as a natural transformation.
  This is the paper's principal open theorem; the obstruction is the doubly-degenerate
  $\chi(M^2)^2$ sector that first appears at $n=5$.
\item \textbf{The novelty gate on $L=-M^4$.} Settle whether the $\SL(4)$ relation
  (Theorem~\ref{thm:m4l}) and the family $L=(-1)^{n-1}M^n$ are new within the
  Bergeron--Falbel--Guilloux / Garoufalidis--Thurston--Zickert frameworks
  \cite{bfg2014,gtz2015}, by a specialist read. Until then the result stays \sGated.
\item \textbf{The principal $\SL(5)$ representation.} Either locate it (against the forced
  degeneracy of Section~\ref{sec:knot}) or prove it cannot exist as an irreducible bundle
  representation --- which would confirm that ``exponent $=$ rank'' is exactly an
  $n\in\{3,4\}$ phenomenon.
\end{enumerate}

\subsection{Method and reproducibility}

The results above were produced under a fixed verification discipline, which we describe
because it is part of what makes a wide-ranging exploration trustworthy. Every figure in
this paper is regenerated by a single deterministic script
(\texttt{figures/gen\_figures.py}); the symbolic and finite-field computations behind
Theorems~\ref{thm:tower}, \ref{thm:m4l} and \ref{thm:complete} are reproducible from the
project's probe code and its committed data artifacts. Five guard-rails, applied
throughout, are worth naming:
\begin{itemize}[leftmargin=1.4em]
\item \emph{Run the control.} A number that reproduces exactly may still be
  interpretively false; vary the causal variable and confirm the effect tracks it.
\item \emph{Filter the reducible locus} before counting character-variety points, and
  test irreducibility with a criterion valid for non-semisimple representations (Burnside,
  not the Schur commutant --- see Section~\ref{sec:knot}).
\item \emph{Generic $\ne$ discriminating.} Check the null case; a feature present
  everywhere distinguishes nothing.
\item \emph{State results twice:} the bare mathematical statement, and --- separately
  tagged, never contaminating the first --- any physical reading
  (Remark~\ref{rem:postulated}).
\item \emph{Check for vacuity.} Before using a target, exhibit a case in which it could
  fail; a property that cannot fail is not a property.
\end{itemize}
These are not stylistic preferences. Each caught a concrete error in the work reported
here: a leading-coefficient bug in the chirality count (Section~\ref{sec:chirality}), a
false irreducibility verdict from the Schur commutant (Section~\ref{sec:knot}), and the
over-statement of a conversational sharpening, demoted to \sPost\
(Remark~\ref{rem:postulated}). The discipline, more than any single theorem, is what we
would ask a reader to trust.

\subsection*{Acknowledgements}
This work was carried out with AI-assisted computation and exposition under the
verification discipline described above; all results are reproducible from the cited code
and data.
